% Glossary of Acronyms
% Cross-disciplinary terminology is defined in Section~\ref{section:language}.

\subsection*{Acronyms}

\begin{longtable}{p{2.2cm} p{12cm}}
A2C/A3C & Advantage Actor-Critic / Asynchronous Advantage Actor-Critic. Policy gradient algorithms that use an advantage function baseline to reduce variance (Section~\ref{section:rl_algorithms}). \\[4pt]
BCQ & Batch-Constrained Q-learning. An offline RL algorithm that restricts the learned policy to actions supported by the behavioral data (Section~\ref{sec:offline_algorithms}). \\[4pt]
BwK & Bandits with Knapsacks. A bandit framework with resource constraints, where an agent must balance exploration and exploitation subject to a budget (Section~\ref{section:bandits}). \\[4pt]
CCP & Conditional Choice Probabilities. Choice probabilities conditional on state, used in structural estimation as an alternative to full-solution methods (Section~\ref{section:rl_econ_models}). \\[4pt]
CFR & Counterfactual Regret Minimization. An iterative algorithm for computing Nash equilibria in extensive-form games by minimizing regret at each information set (Section~\ref{section:rl_games}). \\[4pt]
CQL & Conservative Q-Learning. An offline RL algorithm that adds a penalty for overestimating Q-values on out-of-distribution actions (Section~\ref{sec:offline_algorithms}). \\[4pt]
DDC & Dynamic Discrete Choice. A class of structural econometric models in which agents make sequential discrete decisions under uncertainty (Section~\ref{section:rl_econ_models}). \\[4pt]
DDPG & Deep Deterministic Policy Gradient. An actor-critic algorithm for continuous action spaces that combines a deterministic policy gradient with a learned Q-function (Section~\ref{section:rl_algorithms}). \\[4pt]
DP & Dynamic Programming. A collection of algorithms that compute optimal policies by solving the Bellman equation, given a known model of the environment (Section~\ref{section:rl_algorithms}). \\[4pt]
DPO & Direct Preference Optimization. A method that optimizes a language model directly on preference data without training a separate reward model (Section~\ref{section:rlhf}). \\[4pt]
DQN & Deep Q-Network. Q-learning with a neural network function approximator, target network, and experience replay (Section~\ref{section:rl_algorithms}). \\[4pt]
ETC & Explore-Then-Commit. A bandit algorithm that explores uniformly for a fixed number of rounds, then commits to the empirically best arm (Section~\ref{section:bandits}). \\[4pt]
FQI/FVI & Fitted Q-Iteration / Fitted Value Iteration. Approximate dynamic programming algorithms that use supervised learning to fit value functions from batch data (Section~\ref{sec:fvi_fqi_theory}). \\[4pt]
GLIE & Greedy in the Limit with Infinite Exploration. A condition on exploration schedules ensuring convergence to the optimal policy (Section~\ref{section:rl_algorithms}). \\[4pt]
GMM & Generalized Method of Moments. An econometric estimation method that matches sample moments to population moment conditions (Section~\ref{section:rl_econ_models}). \\[4pt]
IPW & Inverse Probability Weighting. A method for correcting distributional mismatch by reweighting observations by the inverse of their selection probability (Section~\ref{section:causal_rl}). \\[4pt]
IQL & Implicit Q-Learning. An offline RL algorithm that avoids querying out-of-distribution actions by using expectile regression on the value function (Section~\ref{sec:offline_algorithms}). \\[4pt]
IRL & Inverse Reinforcement Learning. The problem of inferring a reward function from observed behavior, assuming the agent acts approximately optimally (Section~\ref{section:rl_algorithms}). \\[4pt]
IV & Instrumental Variables. An econometric technique for identifying causal effects in the presence of endogeneity by exploiting exogenous variation (Section~\ref{section:causal_rl}). \\[4pt]
KL & Kullback-Leibler divergence. A measure of how one probability distribution diverges from a reference distribution, used in trust-region and regularization methods (Section~\ref{section:rl_algorithms}). \\[4pt]
LLM & Large Language Model. A neural network trained on large text corpora, the primary object of alignment via RLHF and DPO (Section~\ref{section:rlhf}). \\[4pt]
LQR & Linear-Quadratic Regulator. An optimal control method that computes the policy for linear dynamics with a quadratic cost function via the Riccati equation (Section~\ref{section:applications}). \\[4pt]
MARL & Multi-Agent Reinforcement Learning. RL in settings with multiple interacting agents, where each agent's optimal policy depends on the others' behavior (Section~\ref{section:rl_games}). \\[4pt]
MC & Monte Carlo. Methods that estimate value functions from complete episode returns rather than bootstrapped estimates (Section~\ref{section:rl_algorithms}). \\[4pt]
MDP & Markov Decision Process. A formal model of sequential decision-making defined by states, actions, transition probabilities, rewards, and a discount factor (Section~\ref{section:rl_algorithms}). \\[4pt]
MLE & Maximum Likelihood Estimation. A statistical method that estimates parameters by maximizing the likelihood of the observed data (Section~\ref{section:rl_econ_models}). \\[4pt]
MPC & Model Predictive Control. A control method that solves a finite-horizon optimization problem at each timestep using an explicit model of the dynamics, then applies only the first action (Section~\ref{section:applications}). \\[4pt]
NE & Nash Equilibrium. A strategy profile in which no player can improve their payoff by unilaterally deviating (Section~\ref{section:rl_games}). \\[4pt]
NFXP & Nested Fixed Point. Rust's algorithm for structural estimation of DDC models that nests value function iteration inside a maximum likelihood loop (Section~\ref{section:rl_econ_models}). \\[4pt]
NPG & Natural Policy Gradient. A policy gradient method that preconditions the gradient by the inverse Fisher information matrix (Section~\ref{section:rl_algorithms}). \\[4pt]
OPE & Off-Policy Evaluation. Estimating the expected return of a target policy using data generated by a different behavioral policy (Section~\ref{section:causal_rl}). \\[4pt]
PEVI & Pessimistic Value Iteration. An offline RL algorithm that subtracts an uncertainty penalty from value estimates to avoid overestimation on unseen state-action pairs (Section~\ref{section:causal_rl}). \\[4pt]
PI & Policy Iteration. A dynamic programming algorithm that alternates between policy evaluation and policy improvement until convergence (Section~\ref{section:rl_algorithms}). \\[4pt]
PID & Proportional-Integral-Derivative controller. A feedback controller that computes a control signal from the weighted sum of the current error, its integral, and its derivative (Section~\ref{section:applications}). \\[4pt]
POMDP & Partially Observable MDP. An MDP in which the agent cannot directly observe the state and must maintain beliefs from noisy observations (Section~\ref{section:causal_rl}). \\[4pt]
PPO & Proximal Policy Optimization. A policy gradient algorithm that constrains updates to a trust region via a clipped surrogate objective (Section~\ref{section:rl_algorithms}). \\[4pt]
RL & Reinforcement Learning. A framework for sequential decision-making in which an agent learns a policy by interacting with an environment and observing rewards (Section~\ref{section:rl_algorithms}). \\[4pt]
RLHF & Reinforcement Learning from Human Feedback. A framework for aligning model behavior with human preferences by training a reward model from pairwise comparisons and optimizing a policy against it (Section~\ref{section:rlhf}). \\[4pt]
SAC & Soft Actor-Critic. An off-policy actor-critic algorithm that maximizes a combined objective of expected return and policy entropy (Section~\ref{section:rl_algorithms}). \\[4pt]
SARSA & State-Action-Reward-State-Action. An on-policy TD control algorithm that updates Q-values using the action actually taken in the next state (Section~\ref{section:rl_algorithms}). \\[4pt]
SFT & Supervised Fine-Tuning. The initial stage of LLM alignment in which the model is trained on curated demonstrations before preference optimization (Section~\ref{section:rlhf}). \\[4pt]
SPE & Subgame Perfect Equilibrium. A refinement of Nash equilibrium requiring that strategies constitute a Nash equilibrium in every subgame (Section~\ref{section:rl_games}). \\[4pt]
TD & Temporal Difference. A class of prediction and control algorithms that update value estimates using bootstrapped targets rather than complete returns (Section~\ref{sec:stochastic_approx}). \\[4pt]
TRPO & Trust Region Policy Optimization. A policy gradient algorithm that constrains each update to lie within a KL-divergence trust region of the previous policy (Section~\ref{section:rl_algorithms}). \\[4pt]
TS & Thompson Sampling. A Bayesian bandit algorithm that selects actions by sampling from the posterior distribution over reward parameters (Section~\ref{section:bandits}). \\[4pt]
UCB & Upper Confidence Bound. A bandit algorithm that selects the arm with the highest sum of empirical mean and an exploration bonus (Section~\ref{section:bandits}). \\[4pt]
VI & Value Iteration. A dynamic programming algorithm that iteratively applies the Bellman optimality operator until the value function converges (Section~\ref{section:rl_algorithms}). \\[4pt]
\end{longtable}
